\section{Introduction}

Les modèles de langage fondés sur l'attention combinent transport contextuel et sélection dépendante du contenu parmi les représentations individuelles des tokens passés. Nous étudions une paramétrisation contextuelle groupée Write--Read--Value (W/R/V) et son compromis qualité--débit face à GQA sous un budget apparié.

Chaque token traverse le même chemin dense. Le mélangeur principal est une attention causale W/R/V à huit têtes de lecture contextuelles et deux têtes partagées d'écriture/valeur. RMSNorm par tête et RoPE précèdent PyTorch SDPA. W/R/V et GQA conservent les mêmes branches feed-forward d'objets lexicaux et de micro-experts déterministes ; les runs W/R/V contrôlés désactivent le résidu lexical d'écriture séparé.

Cette architecture provient d'une étude systématique de WVR à état fixe. Ces variantes sont compatibles avec une limitation d'une matrice comprimée unique pour la sélection d'occurrences. Nous les conservons comme ablations de chemin négatives plutôt que comme architecture finale.

Nos contributions sont empiriques :
\begin{itemize}
    \item une attention contextuelle W/R/V avec têtes de lecture/écriture groupées, normalisation RMS par tête, RoPE et SDPA configuré pour autoriser FlashAttention ;
    \item une comparaison H200 contrôlée à trois seeds contre GQA, à 100~007~936 tokens par run ;
    \item une loss moyenne W/R/V de $4,684432$ contre $4,703035$ pour GQA, W/R/V étant inférieur sur les trois seeds avec 2,86\% de débit en moins ;
    \item des ablations de chemin qui motivent, sans les isoler, des hypothèses sur l'adressage comprimé ou lexical ;
    \item un rapport fondé sur les artefacts bruts qui sépare observations terminées et revendications architecturales appariées en optimisation.
\end{itemize}

Le résultat apparié est cohérent mais faible : avec trois seeds, son intervalle à 95\% contient zéro. Nous ne revendiquons donc ni significativité statistique ni supériorité universelle. La contribution est une observation appariée entièrement archivée et un chemin exécutable à tester à plus grande échelle.
